%% Adapted for Career Pathfinder from the job-application-assistant workflow.
%% Licensed under the MIT License; see README.md (Credits and licence).
%%
%% CV template - [YOUR_NAME]
%% Comprehensive overview of all competencies, experience, and achievements.
%% Base CVs (cvs/base_profile_N_cv.tex) and tailored CVs (cvs/<role-key>_cv.tex) are
%% copied from this file; see .claude/skills/career-cv-template/SKILL.md.
%% Personal details come from personal.tex - never type them into this file.
%%
%% Compile with: lualatex (python3 scripts/build_document.py workspace/cvs/<file>.tex does this)
%% (pdflatex often fails on modern MiKTeX with fontawesome5 font-expansion errors.)

\documentclass[11pt,a4paper,sans]{moderncv}
\moderncvstyle{banking}
\moderncvcolor{blue}

% Force the name and section headings to render in moderncv blue (color1).
% Default banking leaves them black: moderncvstylebanking.sty's \colorlet
% copies (not aliases) the pre-scheme accent colour, so the name colours are
% frozen before \moderncvcolor runs. Re-let them after. \namefont is the hook
% every name-style macro routes through, so this also works on moderncv 2.3.1
% (Debian/Ubuntu apt), which has no \firstnamestyle/\lastnamestyle at all.
\renewcommand*{\namefont}{\fontsize{34}{36}\bfseries\upshape}
\colorlet{firstnamecolor}{color1}
\colorlet{lastnamecolor}{color1}
\colorlet{namecolor}{color1}
\renewcommand*{\sectionstyle}[1]{{\sectionfont\color{color1}#1}}

\usepackage[utf8]{inputenc}
% pdflatex fallback only (the documented engine is lualatex, which skips this
% branch). Without T1 font encoding pdflatex builds accented letters with
% \accent, and the PDF text layer stores them decomposed - `e` + U+0300 rather
% than U+00E8 - so an ATS keyword match on "Genève" fails while the page looks
% right. moderncv 2.5 loads T1 itself under pdflatex; 2.3.1 (Debian/Ubuntu apt)
% does not. \ifpdftex comes from iftex, which every moderncv version loads.
\ifpdftex\usepackage[T1]{fontenc}\fi

% Personal data lives in workspace/personal.tex (created at intake with placeholders;
% the user fills in the real details; Claude never opens it again).
% personal.tex - the header details used by every CV and cover letter.
%
% This copy (templates/personal.tex) is the FORMAT only. At your first run,
% `python3 scripts/setup_tools.py --create-personal` copies it to workspace/personal.tex.
%
% Then open workspace/personal.tex in any text editor and replace each placeholder with
% EITHER a placeholder of the same length as your real detail (e.g. Xxxxxx for Ramesh; keeps
% your details off the AI, and you fill in the real ones at the end, see README.md) OR your
% real detail. Save. The agent never opens this file; the build script copies it into a
% temporary folder. .gitignore keeps it off GitHub.
%
% Special characters in your details need a backslash: & -> \&, % -> \%, _ -> \_, # -> \#

\newcommand{\myfirstname}{XXXXX}
\newcommand{\mylastname}{XXXXXXX}
\newcommand{\myphone}{+XX XXXX XXXXXX}
\newcommand{\myemail}{xxxxxx.xxxxx@xxxxx.xxx}
% Full portfolio link, and the short address printed as the link's visible text
\newcommand{\myportfolio}{https://xxxxxxxx.xxx}
\newcommand{\myportfoliotext}{xxxxxxxx.xxx}
% Address on the CV (optional). Leave the braces empty for no address line.
\newcommand{\myaddress}{}


% moderncv loads hyperref itself in an \AtEndPreamble hook, so \hypersetup
% must go in an \AtEndPreamble of our own: on moderncv < 2.4 a top-level
% \usepackage{hyperref} clashes with the class's own
% \RequirePackage[unicode]{hyperref}. From 2.4.0 the class passes its options
% through \PassOptionsToPackage instead, which is what removes that clash.
\AtEndPreamble{\hypersetup{
    colorlinks=true,
    linkcolor=blue,
    filecolor=magenta,
    urlcolor=blue,
    pdftitle={\myfirstname\ \mylastname\ - CV},
    % Keep pdfpagemode=UseNone: this block runs after moderncv's own
    % \AtEndPreamble (moderncv.cls sets pdfpagemode there), so a FullScreen
    % value here would win and open every CV in fullscreen presentation mode.
    pdfpagemode=UseNone,
}}
\usepackage[scale=0.80]{geometry}
% \needspace{5\baselineskip} before a \cventry keeps an entry's title on the same
% page as its bullets (see "Compile-and-inspect loop" in .claude/skills/career-cv-template/SKILL.md).
\usepackage{needspace}
% \ifdefempty, used for the optional address line below.
\usepackage{etoolbox}

% personal data
\name{\myfirstname}{\mylastname}
% No address to list? Leave \myaddress empty in personal.tex and this line is
% skipped. \address{}{}{} fails with "There's no line here to end" on every
% moderncv version.
\ifdefempty{\myaddress}{}{\address{\myaddress}{}{}}
\phone[mobile]{\myphone}
\email{\myemail}
% The portfolio address is printed as the link text, so it is in the text layer
% (an ATS cannot read a URL that only exists behind link text).
\extrainfo{\href{\myportfolio}{\myportfoliotext}}

\begin{document}

\makecvtitle

% ============================================================
%     PROFILE STATEMENT
% ============================================================

\vspace{6pt}
\small{[Write a 3-4 line profile statement tailored to the role you're applying for. Focus on what makes you uniquely qualified. Example: "Data scientist with a PhD in [field] and X years of industry experience. Combines deep domain expertise in [domain] with strong applied machine learning and Python development skills. Experienced in developing end-to-end ML pipelines, from data ingestion to stakeholder-facing dashboards."]}

% ============================================================
%     CORE COMPETENCIES
% ============================================================

\section{Core Competencies}
\vspace{1pt}
\begin{itemize}

\item \textbf{[Skill Category 1]}: [Specific skills, frameworks, tools. Be concrete.]

\item \textbf{[Skill Category 2]}: [Specific skills, frameworks, tools.]

\item \textbf{[Skill Category 3]}: [Specific skills, frameworks, tools.]

\item \textbf{[Skill Category 4]}: [Domain expertise, methods, approaches.]

\item \textbf{[Skill Category 5]}: [Tools, platforms, software.]

\end{itemize}

% ============================================================
%     PROFESSIONAL EXPERIENCE
% ============================================================

\section{Professional Experience}
\vspace{3pt}
\begin{itemize}

% --- Most Recent Role ---
\item{\cventry{[YYYY-Present]}{[Job Title]}{[Company]}{[City, Country]}{}{\vspace{1pt}
\begin{itemize}
    \item {[Achievement or responsibility 1 - be specific, use the user's real numbers where they exist]}
    \item {[Achievement or responsibility 2]}
    \item {[Achievement or responsibility 3]}
    \item {[Achievement or responsibility 4]}
\end{itemize}}}

\vspace{3pt}

% --- Previous Role ---
\item{\cventry{[YYYY-YYYY]}{[Job Title]}{[Company]}{[City, Country]}{}{\vspace{1pt}
\begin{itemize}
    \item {[Achievement or responsibility 1]}
    \item {[Achievement or responsibility 2]}
    \item {[Achievement or responsibility 3]}
\end{itemize}}}

\vspace{3pt}

% --- Earlier Role ---
\item{\cventry{[YYYY-YYYY]}{[Job Title]}{[Company]}{[City, Country]}{}{\vspace{1pt}
\begin{itemize}
    \item {[Achievement or responsibility 1]}
    \item {[Achievement or responsibility 2]}
\end{itemize}}}

\end{itemize}

% ============================================================
%     EDUCATION
% ============================================================
% Education completed more than 6 years ago: leave the date field empty, e.g.
% \cventry{}{[Degree] in [Field]}{...}. This is the user's decision; don't add dates back.
% In progress: say so inside the entry ("In progress, expected [Month Year].").

\section{Education}
\vspace{1pt}
\begin{itemize}

\item{\cventry{[YYYY-YYYY]}{[Degree] in [Field]}{[Institution]}{[City, Country]}{}{\vspace{1pt}
Thesis: ``[Thesis Title].'' [Brief description of research focus.]
}}

\vspace{3pt}

\item{\cventry{}{[Degree] in [Field]}{[Institution]}{[City, Country]}{}{\vspace{1pt}
[Brief description or key topics.]
}}

\end{itemize}

% ============================================================
%     PUBLICATIONS (optional)
% ============================================================

\section{Publications}
\vspace{1pt}
\begin{itemize}
\item {[Author(s)] ([Year]). [Title]. [Journal/Conference]. \href{[DOI_URL]}{DOI link}}
\end{itemize}

% ============================================================
%     HONORS AND AWARDS (optional)
% ============================================================

\section{Honors and Awards}
\vspace{1pt}
\begin{itemize}
\item{\textbf{[Award Name]}, [Event/Organization] ([Year]).}
\end{itemize}

\end{document}
